\section{Towers, exactly one gauge factor, and sharp trace limitations}
\textit{Supports: G331, G345, n1, n2, n3, n5, n6, n7.}

For a finite tower $R\to B\to F$, top generic degrees satisfy
\[
d_P(F)=\sum_{Q\mid P}[K_Q:K_P]h_Q,
\]
where $h_Q$ is the upper-stage generic degree and every field degree is full. Localizing the finite reduced algebras at $P$ makes them products of fields, proving this identity by vector-space dimension. Closed residue degrees also multiply. Applied to the actual cycles and to (18), this proves the unnormalized ALL-pulled product tower rule, without constant degree or equidimensionality.

Normal-domain traces compose by a single splitting field: group top embeddings by their restriction to the intermediate field. Every intermediate embedding extends; field and separable degrees multiply, hence generic inseparable multiplicities multiply too. One chosen supported valuation lift gives the same Dirac sum, and choice independence makes the equality valid on all supports. For the full conductor trace (24), composition requires BOTH intermediate admissibilities: $H$ must descend after its trace to $B$, and that result must be admissible over $R$. To check the normalization identity, a minimal top source component dominating a minimal $R$ component contracts to a unique minimal $B$ component, by integral incomparability; conversely each such dominating chain is active for the composite. The integral closures of $R/P$ and of its finite domain intermediate algebra in the SAME top field coincide: an element integral over the first is integral over the second, and integrality over the second implies integrality over the first by transitivity. Thus the normal branch diagrams use the same field rings. Grouping embeddings proves the identity there on every supported section, including the lower branch centers. The two assumed FULL descents give equality downstairs by uniqueness, and imply direct admissibility. Direct admissibility alone supplies neither intermediate descent.

Finite-flat determinant traces compose everywhere. At a base support, first split the Artinian fibre algebra of $B$. On a local factor of length $l$ with split residue field, filter by a composition series. The upper algebra is flat over that factor, so tensoring each layer is exact and each layer is its residue-fibre algebra. Splitting the latter into factors of lengths $m_j$ gives total lengths $lm_j$. The composed positive kernels therefore have exactly the direct fibre weights, including all nonreduced special-fibre and inseparable lengths. This proves the tower on every support and on extended nonpositive $D$ functions by finite sums. Determinant norm transitivity also follows generically from field norms and then as a global family polynomial: the reduced base injects into its product of generic fields, so equality of coefficients there is equality in the base. Constant ranks multiply; normalized constant-rank averages consequently compose. Normalized varying-rank averages do not.

The all-field varying-rank example is $A=k[[t]]$, $B=A\times A$, and $F=A\times k[[s]]$, where $t=s^2$ in the second component. The Laurent series parity decomposition proves $[k((s)):k((s^2))]=2$ even in characteristic two. On synchronized sections, direct trace weights are $(1,2)$, total degree three; iterated componentwise normalized averages have weights $(1/2,1/2)$. For $H=(0,-1)$, direct average is $-2/3$, iterated average $-1/2$, and both unnormalized traces are $-2$. On a FINITE branch link, equality of normalized kernels in a branch-surjective tower requires the next-stage generic degree to be constant across each intermediate fibre: comparison of its positive point weights proves necessity, while pulling that common degree out proves sufficiency.

There cannot be a theorem with several arbitrary independently traced source slots. Take the smooth three-dimensional $R=k[x,y,z]_{(x,y,z)}$ and $B=R\times R$, diagonal finite flat of rank two. The actual $C_B$ potentials $H=(-1,0)$ and $V=(0,-1)$ have zero product on both components, since each has a zero input. But $\tau H=\tau V=-1$, and $\mu_R(-1,-1)$ has mass one. Separate normalized traces give base mass $1/4$, still nonzero. This exact counterexample fixes the ONE-slot scope of (23)--(24).

For another base primary ideal $\mathfrak b$, let $t(v)=v_{\mathfrak a}(\mathfrak b)>0$ and $T(v)=v/t(v)$. This is a homeomorphism and $t,t^{-1}$ are bounded continuous. Transport every degree-one function AND obstacle by
\[
u_{\mathfrak b}(Tv)=u_{\mathfrak a}(v)/t(v),\qquad
\phi_{\mathfrak b}(Tv)=\phi_{\mathfrak a}(v)/t(v).
\]
Competitors correspond by G331. Multiplication by a positive continuous factor commutes with USC regularization of nonpositive functions: at finite upper limit use a uniform local bound and continuity of the factor; at bottom, any finite threshold is eventually below a sufficiently negative bound times a locally positive lower bound. This neighborhood argument proves the extended case. Hence both raw and SPECIFIED regularized envelopes obey the same cocycle, including their positive domains. Kernels, fibre extrema and full conductor conditions commute because every lift has the same base gauge. For products G345 gives exactly
\[
\mu_{\mathfrak b}(\boldsymbol u_{\mathfrak b})
=T_*[t\,\mu_{\mathfrak a}(\boldsymbol u_{\mathfrak a})]. \tag{25}
\]
For ideal moments the factor cancels the test rescaling; atomically it changes $\operatorname{ord}_Z(\mathfrak a)$ to $\operatorname{ord}_Z(\mathfrak b)$. Thus there is ONE measure factor, not $t^{n-1}$; all finite-map identities commute with it.

For a separate source primary gauge $\mathfrak c$, put $\lambda(w)=w(\mathfrak aB)$ on $X_{B,\mathfrak c}$ and $\rho_{\mathfrak c}(w)=w|_R/\lambda(w)$. The synchronized source point is $w/\lambda(w)$. A synchronized potential $H_{\mathfrak aB}$ corresponds to
$H_{\mathfrak c}(w)=\lambda(w)H_{\mathfrak aB}(w/\lambda(w))$;
a pulled-back input is $\lambda(w)u(\rho_{\mathfrak c}w)$. Therefore
\[
\rho_{\mathfrak c,*}[\lambda\,\mu_{B,\mathfrak c}(\boldsymbol H_{\mathfrak c})]
=\rho_{\mathfrak aB,*}\mu_{B,\mathfrak aB}(\boldsymbol H_{\mathfrak aB}). \tag{26}
\]
The arbitrary traced slot must be evaluated as $H_{\mathfrak c}(w)/\lambda(w)$ on the corresponding synchronized fibre BEFORE trace. This is the precise separate-source rule; tracing the unconverted slot is incorrect. Nonprimary gauges give only the finite positive partial domain, which need not be compact, and are outside these whole-section envelope statements.
